% !TEX program = lualatex
\documentclass[a4paper,11pt]{ltjsarticle}

\usepackage{amsmath,amssymb,bm}
\usepackage{booktabs}
\usepackage{geometry}
\usepackage{tikz}
\usepackage{enumitem}
\usepackage{hyperref}
\geometry{margin=25mm}
\usetikzlibrary{arrows.meta,positioning,shapes.geometric,fit,calc}

\title{提案手法ver2\_b\\archive評価とpopulation状態観測を分離した\\状態依存BO制御型多目的GP}
\author{}
\date{2026年5月20日}

\newcommand{\clip}{\mathrm{clip}}
\newcommand{\HV}{\mathrm{HV}}
\newcommand{\EI}{\mathrm{EI}}
\newcommand{\ND}{\mathrm{ND}}

\begin{document}
\maketitle

\section{ver2\_b の位置づけ}

本研究の提案手法は，多目的遺伝的プログラミング
(Multi-objective Genetic Programming; MOGP) の世代状態を観測し，
交叉率 \(p_c\)，突然変異率 \(p_m\)，次回制御更新までの世代数 \(k\) を，
文脈付きベイズ最適化 (Bayesian Optimization; BO) により逐次決定する
状態依存・閉ループ制御手法である．

ver2\_b では，提案法の BO 制御則そのものは変更しない．
一方で，多目的 GP の成果評価と探索状態観測を明確に分離するため，
現在集団 \(P_g\)，子個体集合 \(Q_g\)，外部 archive \(A_g\) を導入する．
本版の要点は次の 3 点である．

\begin{enumerate}[leftmargin=3em]
  \item HV，報酬の進捗項，最終 Pareto front は外部 archive \(A_g\) に基づいて評価する．
  \item 多様性，平均木サイズなどの探索状態は現在集団 \(P_g\) に基づいて観測する．
  \item archive の構造重複判定として，トポロジーのみのキーと，トポロジーにノード値を含めるキーを切替可能にする．
\end{enumerate}

これにより，探索中に一度発見された有用な非劣解を失わずに評価しつつ，
BO には現在の探索状態を反映した文脈を与えることができる．

\section{閉ループ制御構造}

提案法では，GP 本体をプラント，BO を制御器とみなす．
制御ステップ \(\ell\) において，BO は現在状態 \(x_\ell\) を入力として受け取り，
\begin{equation}
u_\ell=(p_{c,\ell},p_{m,\ell},k_\ell)
\end{equation}
を出力する．ここで \(k_\ell\) は次回 BO 更新までの世代数である．

GP の世代番号を \(g\)，BO の制御更新ステップを \(\ell\) とすると，
\begin{equation}
g_{\ell+1}=g_\ell+k_\ell
\end{equation}
である．制御区間中は \(p_{c,\ell}\) と \(p_{m,\ell}\) を保持し，
GP を \(k_\ell\) 世代だけ進化させる．

\begin{figure}[htbp]
\centering
\resizebox{0.96\linewidth}{!}{%
\begin{tikzpicture}[
  node distance=9mm and 14mm,
  block/.style={rectangle, rounded corners=2pt, draw, align=center, minimum width=39mm, minimum height=9mm, fill=blue!6},
  controller/.style={rectangle, rounded corners=2pt, draw, align=center, minimum width=39mm, minimum height=10mm, fill=orange!10},
  plant/.style={rectangle, rounded corners=2pt, draw, align=center, minimum width=42mm, minimum height=11mm, fill=green!9},
  data/.style={rectangle, rounded corners=2pt, draw, align=center, minimum width=42mm, minimum height=9mm, fill=purple!7},
  arrow/.style={-Latex, thick},
  labelbox/.style={fill=white, inner sep=1.5pt}
]
\node[block] (obs) {状態観測\\ archive HV, population diversity};
\node[controller, right=29mm of obs] (bo) {文脈付きBO制御器\\ GP surrogate + EI};
\node[data, below=12mm of bo] (act) {制御入力\\ \(u_\ell=(p_c,p_m,k)\)};
\node[plant, below=12mm of obs] (plant) {MOGPプラント\\ \(P_g \rightarrow Q_g\)};
\node[block, below=12mm of plant] (archive) {archive更新\\ \(A_{g+1}=\ND(A_g\cup P_g\cup Q_g)\)};
\node[data, right=29mm of archive] (reward) {報酬・履歴更新\\ \((x_\ell,u_\ell,r_\ell)\)};

\draw[arrow] (obs.east) -- node[labelbox, above] {文脈 \(x_\ell\)} (bo.west);
\draw[arrow] (bo.south) -- node[labelbox, right] {EI最大化} (act.north);
\draw[arrow] (act.west) -| node[labelbox, near start, below] {操作入力} (plant.east);
\draw[arrow] (plant.south) -- node[labelbox, left] {子個体 \(Q_g\)} (archive.north);
\draw[arrow] (archive.east) -- node[labelbox, above] {区間統計} (reward.west);
\draw[arrow] (reward.east) -- ++(12mm,0) |- node[labelbox, pos=0.35, right] {BO更新} (bo.east);
\draw[arrow] (archive.west) -- ++(-13mm,0) |- node[labelbox, near end, left] {次状態} (obs.west);
\end{tikzpicture}
}
\caption{ver2\_b における archive 評価と population 状態観測を分離した閉ループ構造}
\label{fig:closed_loop_v2b}
\end{figure}

\section{現在集団，子個体集合，外部 archive}

世代 \(g\) の現在集団を
\begin{equation}
P_g=\{T_1,\dots,T_N\}
\end{equation}
とする．交叉，突然変異，複製により生成され，評価された子個体集合を
\begin{equation}
Q_g=\{T'_1,\dots,T'_N\}
\end{equation}
とする．

NSGA-II 型の生存選択では，
\begin{equation}
P_{g+1}
=
\mathrm{Select}_{\mathrm{NSGA-II}}(P_g\cup Q_g)
\end{equation}
により次世代集団を決める．
一方，外部 archive は進化操作には戻さず，成果評価と報酬計算のために保持する．
archive 更新は
\begin{equation}
A_{g+1}
=
\ND(A_g\cup P_g\cup Q_g)
\end{equation}
を基本とする．
ここで \(\ND(\cdot)\) は非劣解抽出を表す．

この設計により，子個体 \(Q_g\) の中に有用な非劣解が現れた場合，
たとえ次世代集団 \(P_{g+1}\) に残らなくても，成果集合 \(A_{g+1}\) には保存できる．
したがって，最終的な Pareto front は，最終世代の集団 \(P_T\) ではなく，
探索全体で得られた archive \(A_T\) から得る．

\section{状態ベクトル}

BO に渡す文脈ベクトルは，ver1 と同じ 6 変数を基本とする．
ただし，HV に関する成分は archive ベース，多様性と平均木サイズは population ベースとして明確に分ける．
\begin{equation}
x_\ell =
\left[
\tau_\ell,\;
\widetilde{\HV}^{A}_\ell,\;
\widetilde{\Delta \HV}^{A,\mathrm{rate}}_\ell,\;
D^{P}_\ell,\;
\widetilde{L}^{P}_\ell,\;
\widetilde{s}^{A}_\ell
\right]^\top .
\end{equation}

\begin{table}[htbp]
\centering
\caption{ver2\_b における状態ベクトルの構成}
\begin{tabular}{lll}
\toprule
記号 & 計算対象 & 意味 \\
\midrule
\(\tau_\ell\) & 時間軸 & 世代進行率 \(g_\ell/T\) \\
\(\widetilde{\HV}^{A}_\ell\) & archive \(A_{g_\ell}\) & 現在までの成果集合の性能 \\
\(\widetilde{\Delta \HV}^{A,\mathrm{rate}}_\ell\) & archive \(A\) & 直近区間の世代あたり改善量 \\
\(D^{P}_\ell\) & population \(P_{g_\ell}\) & 現在集団の構造多様性 \\
\(\widetilde{L}^{P}_\ell\) & population \(P_{g_\ell}\) & 現在集団の平均木サイズ \\
\(\widetilde{s}^{A}_\ell\) & archive \(A\) & archive HV が改善しない停滞長 \\
\bottomrule
\end{tabular}
\end{table}

この分担により，BO は「これまでの成果がどこまで進んだか」と
「現在の探索集団がどれだけ多様性を保っているか」を同時に観測できる．

\section{報酬関数}

報酬は，進捗，多様性維持，制御コストの 3 要素で構成する．
進捗項には archive HV の世代あたり改善量を用いる．
\begin{equation}
\Delta \HV^{A,\mathrm{rate}}_\ell
=
\frac{
\HV(A_{g_{\ell+1}})-\HV(A_{g_\ell})
}
{k_\ell}.
\end{equation}
正規化した進捗項は
\begin{equation}
\widetilde{\Delta \HV}^{A,\mathrm{rate}}_\ell
=
\clip
\left(
\frac{\Delta \HV^{A,\mathrm{rate}}_\ell}{\delta_{\HV}^{\mathrm{rate}}},
-1,1
\right)
\end{equation}
とする．

多様性項は現在集団に基づく区間平均多様性である．
\begin{equation}
\bar{D}^{P}_\ell
=
\frac{1}{k_\ell}
\sum_{j=1}^{k_\ell}
D(P_{g_\ell+j})
\end{equation}
を用いる．目標多様性 \(D^\ast(\tau)\) に対する整合度を
\begin{equation}
S_D(\bar{D}^{P}_\ell,\tau_{\ell+1})
=
\exp
\left[
-
\left(
\frac{\bar{D}^{P}_\ell-D^\ast(\tau_{\ell+1})}{\sigma_D}
\right)^2
\right]
\end{equation}
で定義する．

制御コストは
\begin{equation}
C_k(k_\ell)
=
\frac{k_{\max}-k_\ell}{k_{\max}-k_{\min}}
\end{equation}
とし，最終報酬を
\begin{equation}
r_\ell
=
w_p
\widetilde{\Delta \HV}^{A,\mathrm{rate}}_\ell
+
w_d
S_D(\bar{D}^{P}_\ell,\tau_{\ell+1})
-
w_c
C_k(k_\ell)
\end{equation}
とする．

\section{archive における重複除去}

archive \(A_g\) は，目的値と構造キーのペアで重複管理する．
個体 \(T\) の目的値ベクトルを
\begin{equation}
\bm{y}(T)=(y_1(T),\dots,y_M(T))
\end{equation}
とし，丸め後の目的値キーを \(\rho(\bm{y}(T))\) とする．
構造キーを \(\kappa(T)\) とすると，archive 上の同一性キーは
\begin{equation}
\psi(T)
=
\left(
\rho(\bm{y}(T)),
\kappa(T)
\right)
\end{equation}
である．

ver2\_b では，構造キーとして次の 2 種類を切替可能にする．

\subsection{b1: トポロジーのみの構造キー}

test\_ver2\_b1 では，ノードラベルと arity のみから構造キーを作る．
個体 \(T\) の prefix 表現を
\begin{equation}
\left[
(a_1,d_1),
(a_2,d_2),
\dots,
(a_n,d_n)
\right]
\end{equation}
とすると，
\begin{equation}
\kappa_{\mathrm{topo}}(T)
=
\left[
(a_1,d_1),
(a_2,d_2),
\dots,
(a_n,d_n)
\right]
\end{equation}
である．ここで \(a_i\) はノード種別，\(d_i\) は子ノード数である．
この場合，同じ series/parallel 構造を持つ個体は，
ばね定数や減衰係数が異なっていても同じ構造として扱われる．

\subsection{b2: トポロジーとノード値の構造キー}

test\_ver2\_b2 では，ノードラベル，arity に加えて，
ノードが持つパラメータ値も構造キーに含める．
各ノード \(v\) のラベルを \(a(v)\)，値を \(\theta(v)\)，子ノード列を
\(\mathrm{ch}(v)\) とすると，再帰的な構造キーを
\begin{equation}
\kappa_{\mathrm{topo+val}}(v)
=
\left(
a(v),
\rho_\theta(\theta(v)),
|\mathrm{ch}(v)|,
\{\kappa_{\mathrm{topo+val}}(c)\mid c\in \mathrm{ch}(v)\}
\right)
\end{equation}
と定義する．ここで \(\rho_\theta\) はノード値の丸め処理である．

構造探索問題では，ばねノードのばね定数，ダンパノードの減衰係数が
\(\theta(v)\) に対応する．記号回帰テンプレートでは，定数ノードの値が
\(\theta(v)\) に対応する．このため b2 では，
同じトポロジーでもパラメータ値が異なる設計を別構造として archive に残せる．

\section{文脈付き BO 制御器}

BO の学習対象は ver1 と同じく
\begin{equation}
r=f(x,p_c,p_m,k)+\varepsilon
\end{equation}
である．\(k\) を離散候補として扱うため，各 \(k\) ごとに
\begin{equation}
\mathcal{D}_k
=
\{(x_i,p_{c,i},p_{m,i},r_i)\mid k_i=k\}
\end{equation}
を構成し，
\begin{equation}
f_k(x,p_c,p_m)\sim\mathcal{GP}(m_k,K_k)
\end{equation}
を学習する．

現在状態 \(x_\ell\) を固定し，各 \(k\) と \((p_c,p_m)\) 候補に対して
\(\EI_k(x_\ell,p_c,p_m)\) を評価する．制御則は
\begin{equation}
u_\ell^\ast
=
\arg\max_{k\in\mathcal{K}(x_\ell)}
\max_{(p_c,p_m)\in\mathcal{U}_k(x_\ell)}
\EI_k(x_\ell,p_c,p_m\mid\mathcal{D}_k)
\end{equation}
である．

\section{アルゴリズム}

ver2\_b の処理は次の通りである．

\begin{enumerate}[leftmargin=3em]
  \item 初期集団 \(P_0\) を生成し，評価する．
  \item 初期 archive \(A_0=\ND(P_0)\) を作る．
  \item 現在の \(A_{g_\ell}\) と \(P_{g_\ell}\) から状態 \(x_\ell\) を観測する．
  \item BO が \(u_\ell=(p_c,p_m,k)\) を選択する．
  \item \(p_c,p_m\) を保持し，GP を \(k_\ell\) 世代実行する．
  \item 各世代で \(Q_g\) を生成・評価し，\(A_{g+1}=\ND(A_g\cup P_g\cup Q_g)\) を更新する．
  \item NSGA-II 型選択により \(P_{g+1}\) を決定する．
  \item 区間終了後，archive HV 改善率と population 多様性から報酬 \(r_\ell\) を計算する．
  \item \((x_\ell,u_\ell,r_\ell)\) を BO の履歴へ追加する．
  \item 終了条件を満たすまで繰り返す．
\end{enumerate}

\section{ver1 からの変更点}

\begin{table}[htbp]
\centering
\caption{ver1 と ver2\_b の差分}
\begin{tabular}{lll}
\toprule
項目 & ver1 & ver2\_b \\
\midrule
HV の対象 & 明示的分離なし & archive \(A_g\) \\
多様性の対象 & 現在集団を想定 & population \(P_g\) と明記 \\
archive 更新 & 詳細未固定 & \(A_g\cup P_g\cup Q_g\) から更新 \\
最終 Pareto front & 非劣解集合 & archive \(A_T\) から取得 \\
構造重複判定 & 目的値・構造ペア & b1/b2 で構造キーを切替 \\
BO 制御則 & \(f(x,p_c,p_m,k)\) & 変更なし \\
\bottomrule
\end{tabular}
\end{table}

\section{まとめ}

ver2\_b では，提案法の中心である
「GP の状態に応じて \(p_c,p_m,k\) を文脈付き BO で制御する」
という構造は維持する．
そのうえで，成果評価を archive，探索状態観測を population に分離し，
HV・報酬・最終 Pareto front の意味を明確化した．

また，archive における構造重複判定を
\(\kappa_{\mathrm{topo}}\) と
\(\kappa_{\mathrm{topo+val}}\)
の 2 条件で切替可能にした．
これにより，構造探索問題において，
同一トポロジーの別パラメータ設計を同一解としてまとめる場合と，
別設計として保持する場合の両方を比較できる．

\end{document}
